\section{Fundamentação e trabalhos relacionados}

\textbf{Predição de alvos.} Bancos como o ChEMBL~\cite{gaulton:12,zdrazil:24} reúnem medições de potência de moléculas contra alvos proteicos, e é com eles que se treinam modelos que sugerem a quais proteínas uma molécula se liga. Uma abordagem clássica compara conjuntos de ligantes por similaridade química e relaciona proteínas pela semelhança do que elas ligam, o que permitiu prever interações novas depois confirmadas em laboratório~\cite{keiser:07}. Este trabalho segue outra linha, mais simples e interpretável: classificar a molécula a partir de atributos numéricos calculados dela.

\textbf{Descritores e SMILES.} A regra dos cinco de Lipinski~\cite{lipinski:01} usa massa molecular, LogP e números de doadores e aceptores de ligação de hidrogênio para julgar se uma molécula tende a ser absorvida. A TPSA~\cite{ertl:00} estima a área polar da molécula somando contribuições de fragmentos e se relaciona a absorção intestinal e à passagem pela barreira hematoencefálica. O SMILES~\cite{weininger:88} descreve a molécula como texto: letras são átomos, dígitos fecham anéis e parênteses abrem ramificações. Contar esses símbolos dá atributos de composição sem analisar a estrutura, ao custo de não enxergar subestruturas, como o grupo hidroxamato.

\textbf{Árvores de decisão.} O CART~\cite{breiman:84} divide o espaço de atributos de forma recursiva, em cada nó escolhendo o limiar de um atributo que mais reduz a impureza (Gini ou entropia). Sem limite de profundidade, a árvore acerta todo o treino; a poda por custo-complexidade, controlada por $\alpha$, corta ramos que pouco reduzem o erro e diminui o número de folhas. A vantagem é a legibilidade das regras. Como referência de desempenho usam-se conjuntos de árvores, como a Random Forest~\cite{breiman:01}, mais precisos e menos interpretáveis.

\textbf{Vazamento e validação.} Vazamento de dados, quando informação do teste chega ao treino, é uma causa documentada de resultados superestimados em ciência baseada em aprendizado de máquina~\cite{kapoor:23}. Em química, a divisão aleatória já é conhecida por ser otimista; Sheridan~\cite{sheridan:13} mostra que separar treino e teste por data estima melhor o desempenho prospectivo. Aqui o vazamento é mais direto: moléculas distintas com os mesmos cinco descritores podem cair em lados opostos da divisão. Por isso o teste agrupa as linhas pela tupla de descritores (Seção~\ref{sec:splits}).
